
\clearpage
% ============================================================================
\section*{v4.10.0-r73 Source-Only: The Common Lepton Offset Is an Angle; Scale-Free Branch Choice; the Value of the Angle Stays Measured; Balance of the Lepton Masses (H213--H227, k426--k441)}
\addcontentsline{toc}{section}{v4.10.0-r73 Source-Only: The Common Lepton Offset Is an Angle; Scale-Free Branch Choice; the Value of the Angle Stays Measured; Balance of the Lepton Masses (H213--H227, k426--k441)}
\label{sec:v4100-r73-lepton-angle-balance}
% ============================================================================

\paragraph{Purpose.}
Layers r63 and the tests H167/H168 left the lepton phase in the following state: under the equal standing of the frames (\(J\)) the Koide
value \(K=\tfrac23\) is fixed, and the Brannen phase of the charged leptons (Im sector, formerly ``charged''; r69) and of the neutrinos
(Re sector) is offset from the octant values \(\pi/12\) and \(\pi/6\) by a common small number \(x\).  On 28 November 2026 fifteen tests
followed (H213--H227, scripts k426--k441, and the exploration k436), each fixed in an advance note before the computation.  This layer
records them, on the author's decision of 28 November 2026.  It adds no new computation.  The outcome in one sentence: the common offset
is a dimensionless rotation angle \(\varphi\), not a scale; the branch of the data is chosen without a scale by the sign of a cubic spectral
invariant, and that sign is set by the author; the value \(\varphi=0.0395773\) is not derived by any of the routes tried and remains the one
measured number of the lepton picture; with it, three numbers carry all six lepton masses, and the absolute neutrino scale is recorded as a
prediction.  The tests H215--H220 belong to a parallel strand (stability of \(\beta_{\rm sym}\)) and are recorded in r73.3.  The tests
H173--H212 (mixing strand, chiral quark strand, bands on the sphere) are documented in the research notes and are not recorded in this
layer.  Earlier layers are not rewritten.

\noindent\textbf{Source-only status.}
This v4.10.0-r73 layer is source-only.  It does not build a PDF, does not assign a new DOI,
and does not authorize a public release.  The v5.5.0 Version DOI
\href{https://doi.org/10.5281/zenodo.24563811}{10.5281/zenodo.24563811} and the Concept DOI
\href{https://doi.org/10.5281/zenodo.19210728}{10.5281/zenodo.19210728} are unchanged by it.

\subsection*{v4.10.0-r73 setting}
\label{subsec:v4100-r73-setting}

Brannen form in each lepton sector \(s\): \(\sqrt{m_k}=\mu_s\bigl(1+a\cos(\delta_s+2\pi k/3)\bigr)\), \(k=0,1,2\), with \(a=\sqrt2\) (\(K=\tfrac23\),
protected under \(J\)).  Phase potential of H167, written for both sectors,
\[
V_s(\delta)=-\alpha'_s\cos2(\delta-\chi_s-\varphi)+q\cos3\delta ,
\qquad
(\alpha'_{\rm Im},\chi_{\rm Im})=(+1,\pi/4),\quad (\alpha'_{\rm Re},\chi_{\rm Re})=(-1,0),
\]
where \(\varphi\) is the common rotation of the mixing angle \(\chi\) (introduced in r73.4).  Measured: \(x=\pi/12-\delta^*=0.0395773\) with
\(\delta^*=0.2222220\) (under \(J\) from \(m_\tau\); \(\sigma\approx7.8\cdot10^{-7}\)); from the neutrino ratio \(R=\Delta m^2_{21}/\Delta m^2_{31}\):
\(x_\nu=0.04336\pm0.00213\).  Pole masses \(m_e,m_\mu,m_\tau=0.51099895,\ 105.6583755,\ 1776.93\pm0.09\) MeV;
\(\Delta m^2_{21}=(7.537\pm0.097)\cdot10^{-5}\) eV\(^2\), \(\Delta m^2_{31}=(2.511\pm0.0205)\cdot10^{-3}\) eV\(^2\).  Tools: Python (numpy, scipy, sympy) for
every main computation, Wolfram Language independently for every test; all common quantities agree.  Qiskit was not used (no circuit).
Circulant sector, tree level, H99 assumed throughout.

\subsection*{v4.10.0-r73 results}
\label{subsec:v4100-r73-results}

\paragraph{Theoremlet v4.10.0-r73.1 (the neutrino and charged offsets agree at \(1.78\sigma\); correction of a research-note figure; H213, k426; holds).}
The research note of H207 (not part of an earlier layer) quoted a \(4.2\sigma\) distance between \(x_\nu\) and \(x\).  The denominator was the fit
uncertainty \(0.00091\) of H175 instead of \(\sigma(x_\nu)=0.00213\) propagated from \(R\).  Correctly the distance is \(1.78\sigma\) (Monte Carlo
\(1.79\sigma\)); the inverted ordering offers no nearer branch.  A common offset of both sectors (H168) is compatible.  \emph{Correction recorded:}
the figure \(4.2\sigma\) is withdrawn.

\paragraph{Theoremlet v4.10.0-r73.2 (pole masses are the scheme of the frame; H214, k427; holds).}
Only for pole masses does \(K=\tfrac23\) hold (\(-0.4\sigma\)); running masses at a common scale miss \(K\) by \(110\sigma\) (\(M_Z\)) to \(229\sigma\).
Running masses each at its own scale are equivalent to pole masses (\(\Delta x=0.002\sigma\)).  The reading of \(x\) at \(M_Z\) (H213c, \(1.22\sigma\))
is therefore not admissible; \(R\) is scheme independent; the distance stays \(1.78\sigma\).  \(x=x_\nu\) would need \(m_e/m_\mu\) larger by
\(22.6\%\), 441 times any admissible scheme shift.

\paragraph{Theoremlet v4.10.0-r73.3 (parallel strand: no simple loop of the frame attracts \(\beta_{\rm sym}\); H215--H220, k428--k433; refuted).}
H215: among 66\,429 words of length \(\le5\) in the operations of the frame (11\,748 attracting fixed points in \((0,1)\)) three have a fixed
point within the window \(\pm6.5\cdot10^{-4}\) of \(\beta_{\rm sym}\), none exactly; no accumulation above controls (refuted).  H216: inverse
loops give a formal excess (240 words against a control median of 46), but 182 of them share the known fixed point \(x_E=0.5493155\) of the
single word \(E\), \(9\cdot10^{-5}\) off \(\beta_{\rm sym}\) (by wording partly; in substance the known fixed point; the cluster analysis is
post hoc).  H217: the search does not discriminate between the Casimir area \(\sqrt{j(j+1)}\) and \(j+\tfrac12\); it prefers the latter as a closed
form.  H218: for the band tilt \(\chi\) and for \(\Lambda\) no loop and no target value can be derived without a new setting.  H219, H220:
inverse searches for the end numbers with the building blocks of the frame (elementary, then structural with a density rule) find no more
than random targets; every formula found is fitted.  The counting postulate of r67 stands; its Casimir area is not explained by these searches.

\paragraph{Theoremlet v4.10.0-r73.4 (the common offset is a rotation angle, not a scale; H221, k434; partly, criterion ``scale free'').}
With universal dimensionless coefficients the Re sector separates the two ways of H167 by which the electron leaves the octant: the
conjugation term coincides with the spin-1 term for \(c=0\) and gives \(x_\nu=0\) (\(+23\sigma\)); the sixth-order term with a sector scale
shifts the Re sector by \(1.59\) times the Im sector (linear \(\sqrt2\)) and misses by \(9.2\sigma\).  A common rotation \(\chi\to\chi+\varphi\) in both
sectors moves both minima \emph{exactly} by \(\varphi\), to \(\pi/12-\varphi\) and \(\pi/6-\varphi\), for every \(\varphi\) (Wolfram: slope 0,
curvature 4, symbolic).  Hence the common offset is an angle; a scale does not explain it.  Open after this step: the value of \(\varphi\), and the
twofold degeneracy of each minimum without the sixth-order term.

\paragraph{Theoremlet v4.10.0-r73.5 (scale-free branch choice by the cubic spectral invariant; H222, k435; holds, with the sign set).}
\[
\prod_{k=0}^{2}\bigl(1+a\cos(\delta+2\pi k/3)\bigr)=\frac{4-3a^2+a^3\cos3\delta}{4},
\qquad
\frac{\det Y}{\mu^3}=-\frac12+\frac{\cos3\delta}{\sqrt2}\quad(a=\sqrt2)
\]
(exact).  \(\cos3\delta\) is invariant under \(\delta\to-\delta\) and \(\delta\to\delta+2\pi/3\) and changes sign under \(\delta\to\delta+\pi/3\); among
\(\cos n\delta\), \(n=1,\dots,6\), only \(\cos3\delta\) and \(\cos6\delta\) are invariant.  After the rotation the branch pairs are \(\pi/12-\varphi\),
\(\pi/4+\varphi\) (Im; \(\cos3\delta=\pm\cos(\pi/4-3\varphi)\)) and \(\pi/6\mp\varphi\) (Re; \(\cos3\delta=\pm\sin3\varphi\)).  The rule ``larger
\(\cos3\delta\)'' (equivalently: larger \(\det Y\)) chooses in \emph{both} sectors, with the \emph{same} sign, the branch of the data, for
\(\varphi=0.01\ldots0.06\).  It is the limit \(q\to0^-\) of the sixth-order term: the scale drops out, the choice remains.  The rejected Re branch
\(\pi/6+x\) gives \(R=0.0533\) (\(+50.9\sigma\)); for no offset \(x'\in[0,0.1]\) is it within \(2\sigma\).  \emph{Set:} the sign \(q<0\).

\paragraph{Theoremlet v4.10.0-r73.6 (the sign of the branch-choice term follows neither from \(J\), nor from stability, nor from a natural parent term; k436, H223, k437; refuted; author decision).}
Exploration k436 (not blind): the square \(N^2\) of the cubic norm, which would give the sign through stability, is forbidden under \(J\);
the \(J\)-allowed order six (\(\varphi^3r^3\cos3\delta\), \(\varphi^3r^3\cos\delta\)) is equally bounded for \((\beta_3,c_4)\) and \((-\beta_3,-c_4)\)
(exact, \(\varphi\to-\varphi\)).  With \(\cos3\delta\) and \(\cos\delta\) together a wedge of \(43.8\%\) of the directions chooses the data branch in
both sectors.  H223: no natural cubic building block (\(\det W\), \(\operatorname{Tr}W^3\), traceless or pure-trace parts) gives, through its
square under \(J\), the sign the data need: \(\beta_3(J)=4z(27x+9y+z)\), \(+4\) for \(\det W\), \(+36\) for \(\operatorname{Tr}W^3\) (wrong branch), 0
otherwise; a binding to the spin-1 term is excluded because \(\alpha'\) changes sign between the sectors.  Post hoc: the data sign comes only
from the \(J\)-forbidden part of \(N^2\).  \emph{Author decision (28 November 2026):} the sign is treated as a setting.

\paragraph{Theoremlet v4.10.0-r73.7 (\(\varphi\) is not a Wigner rotation of the frame boosts; H224, k438; no hit).}
With the frame rapidities \(\varepsilon=\operatorname{artanh}\beta_{\rm sym}=0.6175296\) and \(2\varepsilon=\operatorname{arcosh}\gamma_L\), the angles
\(\pi/2,\pi/4,2\pi/3,\pi/3\) and divisors \(1,2,3,6\) (48 candidates fixed in advance), \(\tan(\omega/2)=\sin\theta/(\coth\frac a2\coth\frac b2+\cos\theta)\)
gives no strict hit; the nearest \(\omega(\varepsilon,\varepsilon,\pi/4)/3=0.03967\) lies \(+0.23\%\), i.e.\ \(118\sigma\), off.  One loose hit below
\(1\%\) at a random expectation of \(0.066\) (\(p=6\%\)): ``partly'' by the wording of the advance note, not counted in substance.  Wolfram
computed \(\omega\) independently from the polar decomposition of composed \(4\times4\) boosts.

\paragraph{Theoremlet v4.10.0-r73.8 (as a modulus, \(\varphi\) settles at \(\pi/8\); exact exchange of the sector potentials; H225, k439; refuted).}
At \(q=0\) each sector has the minimum \(-1\) for every \(\varphi\): \(\varphi\) is a flat direction.  If \(\varphi\) minimises the common vacuum
energy \(E(\varphi)=\sum_s\min_\delta V_s\), then to leading order \(\tan3\varphi^*=1+\sqrt2/w\) (\(w\) the weight Im\,:\,Re), so \(\varphi^*=\pi/8\) for
\(w=1\) and \(\varphi^*\in(\pi/12,\pi/6)\) for every \(w>0\) (not blind).  Exact minimisation for \(q/\alpha'=-0.001\ldots-0.3\) gives \(\varphi^*=\pi/8\)
exactly for every \(q\); the sixth-order term alone gives \(\pi/8\) or \(7\pi/24\); at \(\varphi=x\) the energy is not stationary.  Post hoc, the
reason:
\[
V_{\rm Re}(-\delta;\pi/4-\varphi)=V_{\rm Im}(\delta;\varphi)
\]
for every \(q\) and every sixth-order coefficient (Wolfram symbolic), so \(E\) is mirror symmetric about \(\pi/8\), where both sectors would have the
same spectrum (\(\delta=\mp\pi/24\)).  The measured \(\varphi\) breaks this exchange.

\paragraph{Theoremlet v4.10.0-r73.9 (no simple exchange-breaking term places the modulus at the measured value; H226, k440; refuted).}
Four terms fixed in advance: (T1) the conjugation term \(c\,\alpha\cos2\delta\) in the Im sector only; (T2) \(q_{\rm Re}=\lambda q\), \(\lambda>0\); (T3)
\(|\alpha'_{\rm Re}|=\kappa\); (T4, control) \(\alpha\cos2\delta\) in both sectors.  T1 makes \(\varphi=x\) stationary for \(\alpha=2.42\,q\) (leading order
exact), with positive curvature, but only as a \emph{local} minimum; the global minimum lies at \(\varphi=-0.70\).  Tuned so that the global
minimum reproduces the charged leptons (\(\alpha=0.120\) at \(q/\alpha'=-0.0852\)), T1 predicts \(x_\nu=0.495\) (\(+212\sigma\)).  T2 would need
\(\lambda=-0.62\) (outside \(\lambda>0\)); T3 does not enter the \(\varphi\)-dependence at leading order.  T4 keeps the exchange (solutions with equal
spectra).  None of 21 frame numbers lies within \(2\%\) of the required ratio.  Post hoc and not counted: under T1 the measured value is a
metastable state.

\paragraph{Theoremlet v4.10.0-r73.10 (balance of the lepton masses; recorded neutrino predictions; H227, k441; partly).}
Six masses, three inputs (\(\varphi,\mu_e,\mu_\nu\)), with \(K=\tfrac23\) and the branch sign set.  Predictions: \(m_\tau\) from \(m_e,m_\mu\)
(Koide) \(1776.969\) MeV (\(+0.43\sigma\); not blind); \(R=0.030839\) (\(+1.80\sigma\); not blind); the absolute neutrino scale.  With \(\mu_\nu\) from
\(\Delta m^2_{31}\), normal ordering, the neutrino roots are \(-0.1958,\ 0.9440,\ 2.2518\) and
\[
m_{1,2,3}=0.379,\ 8.808,\ 50.11\ \text{meV},\qquad
\textstyle\sum m_\nu=59.30\pm0.24\ \text{meV},\qquad
m_\beta=8.87\ \text{meV},\qquad
m_{\beta\beta}\in[1.29,4.00]\ \text{meV}
\]
(\(m_{\beta\beta}\) over free Majorana phases; \(3.49\) meV with the Brannen sign pattern read as a relative sign, a reading).  All values lie below
the current bounds (KATRIN, DESI DR2 with CMB, KamLAND-Zen).  A measured \(\sum m_\nu<58\) meV or an established inverted ordering would refute the
picture.  Exploration (not blind): the scale ratio \(M_\nu/M_e=\mu_\nu^2/\mu_e^2=3.149\cdot10^{-11}\) is no integer power of \(\beta_{\rm sym}\),
\(\gamma_M\), \(1/\gamma_L\) or \(e^{-2\varepsilon}\) (exponents \(40.38,16.82,38.81,19.58\)), and neither \(M_e^2/M_\nu=1.0\cdot10^{10}\) GeV nor
\(v^2/M_\nu=6.1\cdot10^{15}\) GeV is within a factor \(1.1\) of \(v\) or the Planck mass.

\subsection*{v4.10.0-r73 balance of settings}
\label{subsec:v4100-r73-balance}

\begin{center}
\small
\begin{tabular}{@{}>{\raggedright\arraybackslash}p{0.34\textwidth}>{\raggedright\arraybackslash}p{0.58\textwidth}@{}}
\toprule
item & status\\
\midrule
equal offsets of Im and Re sector & compatible (\(1.78\sigma\), pole masses; r73.1, r73.2)\\
offset as a rotation \(\chi\to\chi+\varphi\) & derived: exact equal shift for every \(\varphi\) (r73.4)\\
scale (sixth order with sector scale), conjugation term & excluded as universal source of the offset (r73.4)\\
branch of the data in both sectors & chosen without scale by \(\cos3\delta\propto\det Y\) (r73.5)\\
sign of the branch-choice term & set by the author; not from \(J\), stability or a natural parent term (r73.6)\\
value \(\varphi=0.0395773\) & measured; not a Wigner rotation (r73.7), not a modulus minimum (r73.8), not fixed by a simple exchange-breaking term (r73.9), number searches without hit\\
scales \(\mu_e,\mu_\nu\) and their ratio & set (measured); ratio not a frame number (r73.10)\\
absolute neutrino scale & prediction recorded (r73.10)\\
\(K=\tfrac23\), Brannen form, H99, circulant sector, tree level & set (inputs), as before\\
\bottomrule
\end{tabular}
\end{center}
The open choice moved from ``a small ratio'' (H167: \(\alpha/\alpha'\) or \(q/\alpha'\)) to ``a dimensionless \(x\)'' (H168), to ``a rotation angle and a
branch'' (r73.4), and to ``a rotation angle and a sign'' (r73.5, r73.6).  Since r73.4 it has not moved: four routes left the value of \(\varphi\)
where it was.  The target \(\varphi\) is itself a measured input; whether it is an initial or boundary condition rather than an equilibrium value
(r73.8, r73.9) is a fork left to the author.

\subsection*{v4.10.0-r73 frames}
\label{subsec:v4100-r73-frames}

\begin{center}
\small
\begin{tabular}{@{}>{\raggedright\arraybackslash}p{0.34\textwidth}>{\raggedright\arraybackslash}p{0.22\textwidth}>{\raggedright\arraybackslash}p{0.36\textwidth}@{}}
\toprule
item & frame & comparison\\
\midrule
r73.1, r73.2, r73.4--r73.6, r73.8--r73.10 (masses, phases, potential) & M & pole masses; \(K\) and \(R\) frame independent in form S\\
r73.7 (Wigner rotation) & R \(\to\) M, R \(\to\) T, M \(\leftrightarrow\) T, declared boosts \(\varepsilon\), \(2\varepsilon\) & rotation angle compared with \(\varphi\) in M\\
r73.3 (\(\beta_{\rm sym}\) loops) & M (frame number) & no comparison across frames\\
P & not used & ---\\
\bottomrule
\end{tabular}
\end{center}
The sector exchange of r73.8 (\(\delta\to-\delta\), \(\varphi\to\pi/4-\varphi\), Im \(\leftrightarrow\) Re) is not the frame exchange \(J\) and not time
reversal \(z\mapsto\bar z\); the three are kept apart (r51, r81).

\subsection*{v4.10.0-r73 register entries}
\label{subsec:v4100-r73-register}

\begin{center}
\small
\begin{tabular}{@{}>{\raggedright\arraybackslash}p{0.40\textwidth}>{\raggedright\arraybackslash}p{0.52\textwidth}@{}}
\toprule
statement & status after v4.10.0-r73\\
\midrule
H213, H214, H222 & hold (H222 with the sign set)\\
H221 & partly; criterion ``scale free'' met\\
H215, H219, H220, H223, H225, H226 & refuted\\
H216 & partly by wording; in substance the known fixed point of \(E\)\\
H217, H218 & not discriminating; not derivable\\
H224 & no strict hit (partly by wording, not counted)\\
H227 & partly; neutrino predictions recorded\\
figure \(4.2\sigma\) of the note H207 & withdrawn, replaced by \(1.78\sigma\)\\
sign of the branch-choice term & setting (author, 28 November 2026)\\
\(\varphi=0.0395773\) & measured; the one measured number of the lepton picture\\
\(\sum m_\nu=59.30\pm0.24\) meV, \(m_\beta=8.87\) meV, \(m_{\beta\beta}\in[1.29,4.00]\) meV & predictions, open\\
OP3, OP4 & open\\
\bottomrule
\end{tabular}
\end{center}

\subsection*{v4.10.0-r73 notes}
\label{subsec:v4100-r73-notes}

{\raggedright
\texttt{notes/\allowbreak research/\allowbreak hilbert\_\allowbreak shift\_\allowbreak felder\_\allowbreak 2026\_\allowbreak 11\_\allowbreak 24/}: folders
\texttt{versatz\_\allowbreak spannung\_\allowbreak h213/}, \texttt{massenschema\_\allowbreak h214/}, \texttt{stabilitaet\_\allowbreak rueckwaerts\_\allowbreak h215/},
\texttt{stabilitaet\_\allowbreak folgetests\_\allowbreak h216\_\allowbreak h218/}, \texttt{endzahlen\_\allowbreak suche\_\allowbreak h219/},
\texttt{endzahlen\_\allowbreak struktur\_\allowbreak h220/}, \texttt{skala\_\allowbreak zwei\_\allowbreak sektoren\_\allowbreak h221/},
\texttt{zweigwahl\_\allowbreak skalenfrei\_\allowbreak h222/}, \texttt{vorzeichen\_\allowbreak kubisch\_\allowbreak k436/},
\texttt{elternterm\_\allowbreak vorzeichen\_\allowbreak h223/}, \texttt{wigner\_\allowbreak drehung\_\allowbreak h224/},
\texttt{modulus\_\allowbreak phi\_\allowbreak h225/}, \texttt{tauschbruch\_\allowbreak phi\_\allowbreak h226/}, \texttt{massen\_\allowbreak bilanz\_\allowbreak h227/}:
advance notes, result notes, scripts \texttt{k426} to \texttt{k441} (each with a Wolfram cross-check \texttt{b}) and their actual outputs.
Advance notes of H215--H220 are in \texttt{notes/\allowbreak research/\allowbreak predictions/}.  Layer notes:
\texttt{notes/\allowbreak research/\allowbreak v4\_\allowbreak 10\_\allowbreak 0\_\allowbreak r73\_\allowbreak lepton\_\allowbreak phase\_\allowbreak bilanz/}.
Tools: Python and Wolfram Language for every test; no Qiskit (no circuit); no HDL/MMA comparison (no finite search or selection procedure
of \texttt{RUN\_TAPE} involved).
\par}

\subsection*{v4.10.0-r73 ledger and guard update}
\label{subsec:v4100-r73-ledger-guard}

\[
R_{\rm chk}^{v4100\text{-}r73}=(1,1,1)
\]
for (status recorded; computations checked in the cited scripts with independent Wolfram cross-checks; comparison with data in r73.1, r73.2,
r73.5, r73.9, r73.10; \(\varphi\), the branch sign and the scales not derived).

The layer spends no guard:
\[
R_{\rm glob}^{(v4100\text{-}r73,+)}=(0,1,1,1,1),
\qquad
\Gamma_{\rm guard}^{v4100\text{-}r73}=(1,1,1,1,1),
\qquad
\beta_4=1/256\ \text{unspent}.
\]
The layer does not imply exact finite \(C5\) zero closure, exact all-readout terminality, \(C6/T2\), a global photon-clock theorem,
or global all-readout closure.  OP1--OP5, C5 and C6/T2 remain open.
